\section{Results}\label{results}

\subsection{A multimodal human MASLD Resource records inherited variation and disease remodeling side by side}\label{a-multimodal-human-masld-resource-records-inherited-variation-and-disease-remodeling-side-by-side}

To compare inherited genetic evidence with disease-state evidence gene by gene, we assembled a multimodal human MASLD Resource. It records trait-associated inherited variation (synchronized to the allele-aware, PSD LD-projected eQTL refit of 430 primary multi-signal genes), multicohort disease remodeling, and donor-aware cell and tissue context (deduplicated to 650,083 unique physical cells across 104 source donors) side by side (Figure 1B). For each result, the Resource lists the source cohort, assay, phenotype, biological unit and set of genes tested, whether the same data entered another analysis or source study, the evidence state, the main remaining alternative explanation, and the next discriminating experiment (Figure 1A). The transcriptomic data comprise 1,257 libraries that passed technical quality control, from nine cohorts (GSE126848, GSE130970, GSE135251, GSE162694, GSE167523, GSE174478, GSE193066, GSE213621 and GSE240729). The pooled disease-versus-control model uses 774 genetic-identity clusters from the five cohorts that include controls (138 control and 636 disease clusters). These clusters are not verified independent donors. The genetic data span prespecified liver GWAS strata from several ancestry groups. We kept direct MASLD and liver-imaging traits separate from the liver enzymes ALT, AST, and GGT, which we treat as proxy phenotypes. Donor-collapsed single-cell data, prespecified gene programs, and liver protein, chromatin-accessibility, and spatial measurements add context in each assay's native units. They are not extra votes in a single target score. Throughout, we define replication by participants or biological donors; cells, spots, sequencing runs, guides, and arrays do not add to the biological sample size.

\begin{mainfigureplate}{1}{608bp}
\noindent\includegraphics[width=608bp]{figures/main/figure1/figure1_JL.pdf}
\end{mainfigureplate}
\papercaption{1}


The two maps overlap no more than expected by chance. Of 3,587 genes evaluated by both after the PSD eQTL refit, 420 were colocalized with a liver eQTL (PP.H4 \textgreater{} 0.5), 209 were canonical disease-state genes (BH q \textless{} 0.05 and \textbar log2FC\textbar{} \textgreater{} 0.5) and 26 were both, against 20.9 expected within expression deciles (Mantel-Haenszel OR 1.321, 95\% CI 0.851 to 2.051; exact stratified \emph{p} = 0.222; Figure S1A). We state no direction from this overlap; the full comparison, the genes outside it and the trait sources of colocalization support (Figure S1B) are in Supplementary Note 3. Figures 2 and 3 therefore partition each map on its own terms: colocalized loci by whether they carry MASLD or liver-fat liability, and disease-state genes by whether their change can be separated from cell mix and placed inside hepatocytes.


The Resource is organized as a graph of source dependencies and an analysis workflow that ends in a decision for each gene (Figure 1A). Figure 2 shows that ALT, liver fat and cirrhosis share a minority of their heritability: a few ALT loci are also associated with liver fat and colocalize with plasma lipid species, a minority of individual ALT/AST loci (26 of 167, 15.6\%) pass per-locus elevation of other hepatocyte plasma proteins while the cross-locus slope is positive (0.333 protein SD per MVP ALT SD with zero sample overlap; 0.435 per Pan-UKBB ALT SD), and TNFSF10 is the worked example; predicted accessibility appears only as an annotation beside measured evidence. Figure 3 defines multicohort, cross-sectional stage-associated remodeling, tests it in 493 held-out livers and finds that of 7,214 replicated stage genes, 1,276 change inside hepatocytes in uncorrected single-nucleus pseudobulk (188 after empirical ambient-RNA subtraction, 0 of 9 controls called). Figure 4 places the single-cell programs in those held-out livers and examines three hepatocyte results: lower GNMT, MAT1A and CYP2C19, a CCL20+ state whose CCL20 RNA is mostly in cells with hepatocyte expression profiles, and its relation to PNPLA3 genotype. Figure 5 places candidate genes and programs in molecular and physical tissue context. Figure 6 organizes the two maps in the MASLD Gene Catalog by disease evidence, direction, hepatocyte expression, compounds and trial results, and names the smallest experiment that could resolve each remaining alternative. Supplementary Figure S7 and its Supplementary Note report what predictive models learn from the Resource. No step requires agreement across all modalities, and no aggregate rank is used as a calibrated measure of biological truth.

Two sets of examples recur through these figures. GNMT, MAT1A, and CYP2C19 connect the bulk RNA, cross-sectional stage, protein and per-hepatocyte views of hepatocyte metabolic remodeling. The ECM/IGFBP7 and ductular-injury/BICC1 programs connect prespecified single-cell programs to their stage effects in held-out bulk livers. We report the colocalization state of each example, but that state does not determine its role in disease remodeling.

Taken together, this design lets a reader trace why a gene was nominated, see what has been tested for it, and choose an experiment that could distinguish the remaining alternatives. None of these steps depends on the outcome of a perturbation screen. A separate prospective perturbation study may use the Resource, but its gene list, guides, phenotypes, hits, and mechanisms are not part of this paper.
